\section{One Thing, Many Names}

We have collected several words: vibe, tone, note, fill, will, and earlier link. It would be easy to come away thinking the theory has six ingredients. It does not. It has one. This short chapter makes sure the vocabulary does not trick you into seeing more stuff than there is, because the strict one-substance claim is the spine of everything.

\subsection{The words name angles, not parts}

There is only the vibe. The other words are the vibe, and what the vibe does, seen from different angles. They are like the words ``author,'' ``speaker,'' and ``neighbor'' for one person: different roles of the same individual, not three different people.

\begin{itemize}
\item The \textbf{vibe} is the thing itself, a unit of experience.
\item The \textbf{tone} is the vibe regarded as a felt state, its charge from the inside.
\item The \textbf{will} is that same state regarded from the outside, what the vibe shows its neighbors.
\item The \textbf{note} is one vibe's act of experiencing another. Not a thing between them, an act.
\item The \textbf{link} is that same act seen as a graph edge, the bare connection, when we are thinking in dots and lines.
\item The \textbf{fill} is that same act regarded as carrying a strength and a sign.
\end{itemize}

Read that list again and notice that nothing on it is a separate substance. Tone and will are one state, inside and outside. Note, link, and fill are one act, described three ways. The only entity is the vibe.

\subsection{Why this matters so much}

It would be easy, and wrong, to imagine the crystal as cells (vibes) plus strings (notes) plus weights (fills) plus labels (tones), four kinds of stuff bolted together. If that were the picture, the theory would not be a monism at all, and its central claim, that everything is one substance, would collapse on the first page.

So hold the discipline. When the walkthrough later says ``the fill of a note steers the vibe,'' do not picture a weight on a string pulling a ball. Picture one vibe experiencing another in a way that has a certain strength and sign. When it says ``a signal travels along the links,'' do not picture pulses running down wires. Picture vibes experiencing their neighbors, and those neighbors their neighbors, the act rippling outward. The wires, weights, and labels are figures of speech for one thing experiencing another.

\subsection{The smallest possible world}

This is what makes the theory so spare, and so ambitious. The entire inventory of reality is: a unit of feeling, and the act of one such unit feeling another. From that, and only that, plus the rule and the growth still to come, the theory means to build space, time, matter, force, gravity, the quantum, life, and mind. If it needed a second ingredient, the project would be less remarkable. It does not. One thing, and what it does.

Keep this chapter in your back pocket. Any time the language starts to feel like it is multiplying entities, come back here and collapse it. There is the vibe, and there is the vibe experiencing the vibe. Everything else is a way of pointing.

\textbf{Takeaway:} tone, will, note, link, and fill are not separate ingredients. They are the one vibe, and its one act of experiencing, seen from different angles. There is only the vibe and what it does, and the whole universe is built from that alone.
